\section{Braiding}

\begin{exercise}[Fractional shot noise]
\textbf{Experiment.}
Shot-noise measurements in a fractional quantum Hall state infer mobile
charges of magnitude \(e/3\), smaller than the electron charge
\cite{de-picciotto}.  Suppose rare tunneling events transfer quasiparticle
charge \(q\) independently across a constriction.

Derive the low-frequency relation between current noise and mean tunneling
current for a Poisson process.  Show how the measured ratio determines
\(|q|\) without assigning a microscopic path to each event.  Explain why the
result supports fractionalized excitations rather than electrons arriving in
a regular correlated train, and state the additional assumptions needed for
that inference.
\end{exercise}

\begin{exercise}[An anyonic exchange phase]
\textbf{Experiment.}
Electronic interferometry in the \(\nu=1/3\) fractional quantum Hall regime
shows phase slips consistent with the statistical phase of Abelian anyons
\cite{nakamura}.  The configuration space of indistinguishable particles in a
plane has braid group \(B_n\) as its fundamental group.

Show that a one-dimensional unitary representation of \(B_n\) assigns the
same phase \(e^{i\theta}\) to every elementary exchange.  Compare this with
the only one-dimensional exchange possibilities for particles in three
spatial dimensions.  Compute the change in interference phase when one
quasiparticle winds around another, carefully distinguishing exchange from a
full winding.  Explain how discrete phase slips can reveal \(\theta\) even
when ordinary electromagnetic flux also contributes to the total phase.
\end{exercise}

\begin{exercise}[Fractional collision statistics]
\textbf{Experiment.}
Two dilute beams of \(e/3\) quasiparticles sent into an electronic beam
splitter produce current correlations intermediate between bosonic bunching
and fermionic antibunching \cite{bartolomei}.  Model the incoming particles by
Abelian exchange phase \(e^{i\theta}\).

List the indistinguishable two-particle alternatives contributing to an
output coincidence.  Show where their relative exchange phase enters the
current cross-correlation.  Explain why the measured dependence on source
current and splitter transmission can determine \(\theta\), and recover
\(\theta=\pi/3\) for the reported data.  State which interaction and
independence assumptions are required before calling the correlation a
statistics measurement.
\end{exercise}

\begin{exercise}[Mach--Zehnder braiding]
\textbf{Experiment.}
A chiral electronic Mach--Zehnder interferometer in the fractional quantum
Hall regime displays a phase dependence consistent with braiding an
interfering \(e/3\) quasiparticle around localized quasiparticles
\cite{kundu}.  Represent the two paths by amplitudes and the winding by a
unitary character of the braid group.

Separate the electromagnetic and statistical contributions to the measured
phase.  Explain why adding one enclosed quasiparticle shifts the interference
pattern by the double-exchange phase.  Derive the predicted shift for
\(\theta=\pi/3\).  Explain how the two-path geometry and independent charge
measurement help rule out a purely Coulombic explanation.
\end{exercise}

\begin{exercise}[Fermionic antibunching]
\textbf{Experiment.}
Two identical fermions incident on a two-input beam splitter show suppressed
coincident occupation of the same output mode, opposite to bosonic bunching.
Represent their joint state in the exterior square of the one-particle mode
space.

Apply the one-particle scattering unitary to the exterior product and compute
the output probabilities.  Explain why equal one-particle states wedge to
zero.  Show how partial distinguishability weakens the suppression, and
identify exchange antisymmetry rather than a repulsive force as the
explanation of the ideal effect.
\end{exercise}

\begin{exercise}[Quasiparticle charge from flux insertion]
\textbf{Experiment.}
Adiabatically increasing flux through a puncture in a fractional Hall fluid
transports charge between its boundaries.  Regard the family of gapped ground
states as a bundle over the flux parameter.

Relate the transported charge to the Hall response and hence to the
experimentally measured filling fraction.  Explain why one flux quantum
creates a localized excitation carrying fractional charge.  Show why the
argument depends on the bulk gap and gauge equivalence after a full flux
cycle, and connect this transported charge with the shot-noise measurement.
\end{exercise}

\begin{exercise}[Coherent anyon bunching]
\textbf{Experiment.}
An electronic interference experiment observes coherent bunching and
dissociation of fractional quantum Hall anyons \cite{ghosh-bunching}.  Model
the paths by amplitudes whose relative phase includes both propagation and
fractional exchange.

Identify the alternatives which end in the same detected charge
configuration and must therefore be added before squaring.  Show how the
exchange phase changes enhancement into suppression as the control phase is
varied.  Explain why simultaneous measurements of currents and correlations
separate coherent anyon processes from independent Poisson events.  State
which loss of path coherence would erase the effect.
\end{exercise}

\begin{exercise}[Stable individual anyons]
\textbf{Experiment.}
A graphene Fabry--Perot interferometer resolves configurations containing a
fixed number of localized anyons and finds them stable for minutes or hours
while the interference phase records their statistics \cite{samuelson}.
Let the localized configuration label a sector of the interferometer's state
space.

Explain why a transition between sectors produces a discrete phase jump
rather than a continuous drift.  Relate the jump to the holonomy of an edge
anyon winding around the localized sector.  Show how long dwell times allow
separate estimates of within-sector interference and sector-changing rates.
Explain what the stability demonstrates and why it does not by itself
establish non-Abelian fusion.
\end{exercise}

\begin{exercise}[Thermal transport at a Hall edge]
\textbf{Experiment.}
Heat carried along a fractional quantum Hall boundary is observed to be
approximately quantized in units set by temperature and fundamental
constants, with the sign recording net propagation direction
\cite{banerjee-heat}.
Associate a chiral contribution to each protected edge sector.

Show why independent sectors add while oppositely directed sectors subtract.
Relate the net coefficient to the boundary theory's chiral invariant.
Explain how the measurement constrains candidate quasiparticle theories even
when they have the same electrical Hall conductance.  Identify equilibration
and heat leakage as assumptions needed before interpreting the measured value
topologically.
\end{exercise}
